\ReportSection{7. AI Assistance Disclosure}
\ReportSubsection{7.1 AI Tools Used}
We used AI to help with repetitive implementation work and to make our project
documents more concise. We wrote the main user stories, feature files, and
scenario outlines, and chose the behavior we wanted to test. We also wrote the
expected failures in \Path{BUG_REPORT.md}. AI helped us shorten those descriptions
while keeping the expected behavior, observed results, and impact clear.

We wrote JavaScript code for the suite and used AI to help with its
implementation, particularly the helpers and the consistency of code written by
different team members. This included organizing shared steps, applying common
request and assertion patterns, and helping the separate stories work together
in one suite. We used Codex to help set up the LaTeX formatting, format our
writing into paragraphs, tables, and lists, and revise the wording.

\ReportSubsection{7.2 Verification of AI Output}
We ran the suite and compared the results with the expected behavior in our
scenarios and bug report. Running with different random seeds checked whether
the cases behaved consistently when their order changed. The restoration checks
and service-stopped run also checked the setup and cleanup behavior. These checks
helped us assess whether the shared code supported the tests we intended to run.

For the written report, we reviewed the text against our test design and
recorded results. Shorter wording still needs to preserve
what we tested, what we expected, and what actually happened. The LaTeX output
was also compiled and checked for layout issues.

\ReportSubsection{7.3 Strengths and Limitations}
AI helped us apply common patterns across similar tests and reduce repeated
code. This was useful when combining work from several people, because the
stories needed to use the same helpers and behave consistently within the suite.
It also helped us organize the results and express our findings more clearly.

However, AI can misunderstand an expected outcome or suggest an assertion that
checks the wrong thing. A test passing only shows that its assertions were met.
It does not establish that those assertions describe the intended behavior.
The same applies to the report: making an explanation shorter is useful only if
it still matches the scenario and its evidence.

\ReportSubsection{7.4 Prompt Examples}
\begin{enumerate}[leftmargin=1.6em,itemsep=6bp,topsep=4bp]
\item ``Implement this text block ... into the LaTeX and break it up into a numbered list form.''
\item ``In my categories test suite I implemented helpers. I know that in this file it does something similar. Can you abstract it to use the same helper?''
\item ``Can you make the BUG\_REPORT markdown more concise?''
\end{enumerate}
